\chapter{GLOSSÁRIO} 
\label{cap:glossario}

\begin{description}[leftmargin=3.5cm, style=nextline]

    \item[\textbf{ADB (\textit{Android Debug Bridge})}]
    Ferramenta de linha de comando oficial do AOSP que opera em modelo cliente-servidor, permitindo a comunicação, envio de instruções e automação de \\textit{scripts} entre a estação de trabalho e o dispositivo Android.
    
    \item[\textbf{AOSP (\textit{Android Open Source Project})}] 
    
    Projeto de código aberto mantido pela Google que disponibiliza a estrutura base e a arquitetura do sistema operacional Android. 
    
    \item[\textbf{API (\textit{Application Programming Interface})}]
    Interface de Programação de Aplicações; conjunto de rotinas e padrões de software que permite a comunicação entre diferentes componentes de sistema. 
    
    \item[\textbf{ART (\textit{Android Runtime})}] 
    
    Ambiente de execução do Android responsável por traduzir o código das aplicações em instruções de máquina executadas pelo processador. 
    
    \item[\textbf{Blacklist}] Arquivo de configuração estruturado (em formato JSON) utilizado para catalogar e especificar os identificadores de pacotes (\textit{package names}) alvos de remoção.

    \item[\textbf{Bloatware}] 
    Softwares e aplicativos pré-instalados por fabricantes ou operadoras que operam em segundo plano e consomem memória e ciclos de processamento desnecessários. 
    \item[\textbf{CLI (\textit{Command Line Interface})}] Interface de linha de comando utilizada para a execução de programas e \textit{scripts} em modo texto no terminal. 
    
    \item[\textbf{CPU (\textit{Central Processing Unit})}] Unidade Central de Processamento; componente de hardware responsável pela execução dos cálculos e instruções do sistema.

    \item[\textbf{CVS (\textit{Computer Vision Syndrome})}] 
    Síndrome da Visão do Computador; condição ocupacional caracterizada por fadiga e desconforto ocular resultantes do uso prolongado de telas e terminais de vídeo. 
    
    \item[\textbf{Daemon}] Processo do sistema operacional que executa de forma contínua em segundo plano sem a necessidade de interatividade direta com o usuário (ex.: \texttt{adbd}). 
    
    \item[\textbf{Debloat}] Processo sistemático de limpeza e desinstalação de pacotes, aplicativos e serviços redundantes de um sistema operacional. 
    
    \item[\textbf{Deep Work}] Trabalho profundo; conceito que se refere ao estado de concentração contínua e redução de estímulos/distrações induzido por dispositivos de propósito único. 
    
    \item[\textbf{Desejabilidade Composta}] Modelo de análise estatística e multicritério utilizado para calcular o ponto de equilíbrio ótimo entre métricas de desempenho (\textit{throughput}) e consumo de energia.
    
    \item[\textbf{Dispositivo Legado (\textit{Legacy Device})}] Hardware perfeitamente funcional que se tornou obsoleto em relação às exigências de softwares e atualizações recentes. 
    
    \item[\textbf{DPI (\textit{Dots Per Inch})}] Densidade de pontos por polegada; métrica utilizada para ajustar a escala e a proporção dos elementos visuais da interface do Android.
    
    \item[\textbf{E-ink / E-reader}] Tecnologia de tela de tinta eletrônica reflexiva e dispositivos leitores de livros digitais projetados para alta ergonomia visual. 
    
    \item[\textbf{E-waste}] Resíduos de equipamentos eletroeletrônicos acumulados devido ao descarte e à obsolescência de hardware. 
    
    \item[\textbf{GMS (\textit{Google Mobile Services})}] Conjunto de aplicativos, telemetrias e serviços proprietários da Google que executam em segundo plano no sistema Android. 
    
    \item[\textbf{HAL (\textit{Hardware Abstraction Layer})}] Camada de Abstração de Hardware que faz a ponte entre o \textit{framework} do Android e os \textit{drivers} do Kernel Linux. 
    
    \item[\textbf{IPC (\\textit{Inter-Process Communication})}] Comunicação entre processos; mecanismo que permite a troca segura de dados e chamadas entre diferentes processos no sistema (como a arquitetura Binder). 
    
    \item[\textbf{JSON (\textit{JavaScript Object Notation})}] Formato leve e estruturado para troca e armazenamento de dados em pares de chave-valor. 
    
    \item[\textbf{Kernel}] Núcleo do sistema operacional (baseado em Linux no Android) responsável pela alocação de memória, escalonamento e comunicação com o hardware. 
    
    \item[\textbf{Launcher}] Aplicativo gerenciador da tela de início e da organização de ícones e interfaces gráficas do Android. 
    
    \item[\textbf{Minimalismo Digital}] Filosofia de otimização do uso de ferramentas digitais para reduzir a fragmentação da atenção e focar em tarefas prioritárias. 
    
    \item[\textbf{RAM (\textit{Random Access Memory})}] Memória de acesso aleatório onde são armazenados os dados temporários dos processos em execução. 
    
    \item[\textbf{Settings.db}] Banco de dados interno do Android responsável por armazenar as configurações globais e comportamentos de sistema. 
    
    \item[\textbf{Shell}] Interpretador de comandos do sistema operacional Unix/Linux responsável por receber e executar comandos de terminal. 
    
    \item[\textbf{Slowdown}] Métrica de degradação de desempenho que avalia o atraso relativo de um processo sob condições de carga do sistema. 
    
    \item[\textbf{Software Aging}] Envelhecimento de software; conceito que explica a perda progressiva de fluidez em sistemas motivada pelo acúmulo de serviços e atualizações pesadas em hardwares antigos.
    
    \item[\textbf{SurfaceFlinger}] Serviço do Android responsável por compor as superfícies visuais e controlar a taxa de atualização enviada ao display. 
    
    \item[\textbf{UI (\textit{User Interface})}] Interface do usuário; conjunto de componentes visuais através dos quais o usuário interage com o sistema. 
    
    \item[\textbf{Underclocking}] Redução proposital da frequência de operação do processador com o objetivo de diminuir o consumo de bateria e a dissipação térmica. 
    
    \end{description}
